\section{Caract\'eristiques de l'environnement}

Le Tableau~\ref{tab:cpu-general} regroupe les informations g\'en\'erales sur
l'architecture processeur visibles depuis WSL.

\begin{table}[!t]
\caption{Informations g\'en\'erales du processeur}
\label{tab:cpu-general}
\centering
\small
\begin{tabular}{p{0.43\columnwidth} p{0.47\columnwidth}}
\toprule
\textbf{Champ} & \textbf{Valeur} \\
\midrule
Architecture & x86\_64 \\
Modes d'op\'eration CPU & 32-bit, 64-bit \\
Tailles d'adressage & 46 bits physiques, 48 bits virtuels \\
Ordre des octets & Little Endian \\
Nombre de CPU logiques & 16 \\
Liste des CPU en ligne & 0-15 \\
Identifiant du fabricant & GenuineIntel \\
Nom du mod\`ele & Intel(R) Core(TM) Ultra 7 155H \\
Famille CPU & 6 \\
Mod\`ele & 170 \\
Threads par coeur & 2 \\
Coeurs par socket & 8 \\
Sockets & 1 \\
Stepping & 4 \\
BogoMIPS & 5990.39 \\
Virtualisation & VT-x \\
Fournisseur de l'hyperviseur & Microsoft \\
Type de virtualisation & full \\
\bottomrule
\end{tabular}
\end{table}

Le Tableau~\ref{tab:cpu-cache} pr\'esente les informations de cache et
d'organisation NUMA disponibles dans la m\^eme sortie \texttt{lscpu}.

\begin{table}[!t]
\caption{Caches et topologie m\'emoire}
\label{tab:cpu-cache}
\centering
\small
\begin{tabular}{p{0.43\columnwidth} p{0.47\columnwidth}}
\toprule
\textbf{Champ} & \textbf{Valeur} \\
\midrule
Cache L1d & 384 KiB (8 instances) \\
Cache L1i & 512 KiB (8 instances) \\
Cache L2 & 16 MiB (8 instances) \\
Cache L3 & 24 MiB (1 instance) \\
Nombre de noeuds NUMA & 1 \\
CPU du noeud NUMA 0 & 0-15 \\
\bottomrule
\end{tabular}
\end{table}

Ces valeurs ont \'et\'e mesur\'ees directement depuis WSL. Elles diff\`erent donc
d'une configuration \`a l'autre selon la machine h\^ote et les ressources
allou\'ees \`a l'environnement virtuel.
